\section{Conclusiones}

La realización de esta práctica permitió comprobar experimentalmente el fenómeno de resonancia eléctrica en un circuito RLC serie. A partir de la variación de la frecuencia de la señal senoidal de entrada, se observó que el comportamiento del circuito cambia dependiendo de la relación entre la reactancia inductiva y la reactancia capacitiva.

\vspace{5pt}

Se identificó que, para frecuencias menores a la frecuencia de resonancia, el circuito presentó un comportamiento predominantemente capacitivo, lo cual se reflejó en un desfase negativo entre la señal de entrada y la señal de salida. En cambio, para frecuencias mayores a la resonancia, el circuito presentó un comportamiento predominantemente inductivo, observándose un desfase positivo.

\vspace{5pt}

En la frecuencia de resonancia, las señales medidas en el osciloscopio se encontraron prácticamente en fase, lo que confirma que la parte reactiva de la impedancia se anuló y que el circuito se comportó principalmente como una carga resistiva. En esta condición, la impedancia fue mínima y la corriente alcanzó su valor máximo, por lo que el voltaje medido sobre la resistencia también presentó un valor elevado.

\vspace{5pt}

Además, se determinaron experimentalmente las frecuencias de corte inferior y superior, correspondientes aproximadamente a desfases de $-45^\circ$ y $45^\circ$. Con estos valores se obtuvo el ancho de banda del circuito y el factor de calidad. El valor calculado de $Q$ indicó que el circuito no fue altamente selectivo, ya que presentó una respuesta en frecuencia relativamente amplia.

\vspace{5pt}

Finalmente, se concluye que las diferencias entre los valores teóricos y experimentales pueden atribuirse a las tolerancias de los componentes, la resistencia interna de la bobina, las pérdidas parásitas del montaje y las limitaciones propias de los instrumentos de medición. A pesar de estas variaciones, el fenómeno de resonancia se observó claramente y pudo analizarse mediante las señales obtenidas en el osciloscopio